% Source: https://artofproblemsolving.com/community/c3h2576321_a_couple_problems
% Stable identity: aops-topic-2576321-q01
There are $\dbinom92=36$ equally likely pairs. Let the chosen marbles be $i<j$. Going clockwise around the circumscribed circle, the three vertices cut it into arcs of
\[g_1=i-1,\qquad g_2=j-i,\qquad g_3=11-j\]
steps of $36^\circ$, where $g_1+g_2+g_3=10$ and each $g_k\ge1$. Each angle of the triangle is an inscribed angle subtending one of these arcs, so it equals $18^\circ\cdot g_k$.

The centre lies inside or on the triangle exactly when no angle is obtuse, i.e. every $g_k\le5$. (If an angle is obtuse, the centre lies beyond the opposite side; if it is right, the centre is the midpoint of the hypotenuse.)

The pairs $(i,j)$ correspond one-to-one to the triples $(g_1,g_2,g_3)$ of positive integers with sum $10$. A triple with some $g_k\ge6$ has its other two parts summing to at most $4$. For a given position $k$, put $g_k=5+e$ with $e\ge1$; then $e$ and the other two parts are positive with sum $5$, which gives $\dbinom42=6$ triples. Only one part can be $\ge6$, so $3\cdot6=18$ triples are bad, and $36-18=18$ are good.

The probability is $\dfrac{18}{36}=\boxed{\dfrac12}$. In the sketch, triangle $1,4,8$ (arcs $3,4,3$, shaded) contains the centre, while triangle $1,3,5$ (arcs $2,2,6$) does not.

\begin{asy}
usepackage("amsmath");
size(6.5cm);
defaultpen(fontsize(10pt));
void deca(pair c,int i,int j,bool shade){
 pair[] V;
 for(int k=0;k<10;++k) V.push(c+dir(90-36*k));
 draw(V[0]--V[1]--V[2]--V[3]--V[4]--V[5]--V[6]--V[7]--V[8]--V[9]--cycle);
 path t=V[0]--V[i-1]--V[j-1]--cycle;
 if(shade) fill(t,gray(0.85)); draw(t,linewidth(0.8));
 for(int k=0;k<10;++k){dot(V[k]);label("$"+string(k+1)+"$",V[k],V[k]-c,fontsize(9pt));}
 dot(c);label("$O$",c,N,fontsize(9pt));}
deca((0,0),4,8,true);
deca((2.9,0),3,5,false);
\end{asy}
